\documentclass[preprint,12pt]{elsarticle}
\usepackage[T1]{fontenc}
\usepackage[utf8]{inputenc}
\usepackage{amsmath,amssymb,amsthm,mathtools,bm,booktabs}
\usepackage{graphicx,hyperref}
\journal{Digital Signal Processing}


\newtheorem{theorem}{Theorem}
\newtheorem{proposition}{Proposition}
\newtheorem{corollary}{Corollary}
\newtheorem{remark}{Remark}
\newcommand{\KL}{D_{\mathrm{KL}}}
\newcommand{\TV}{\mathrm{TV}}
\newcommand{\Exp}{\operatorname{Exp}}
\newcommand{\Vol}{\operatorname{Vol}}
\newcommand{\R}{\mathbb R}
\newcommand{\E}{\mathbb E}
\newcommand{\Var}{\operatorname{Var}}
\newcommand{\cN}{\mathcal N}
\newcommand{\sgn}{\operatorname{sgn}}
\newcommand{\PhiN}{\Phi}


\begin{document}
\begin{frontmatter}
\title{Fisher--Rao-Guided Generative Steganography:
Geodesic Embedding and Detection Analysis}
\author[uy1]{St\'ephane Ga\"el R. EKODECK}
\author[uy1]{Prosper Rosaire MAMA ASSANDJE}
\author[uy1]{Serge Alain EB\'EL\'E}
\author[uy1]{St\'ephane Willy MOSSEBO TCHEUNTEU}
\author[uy1]{Arthur Ulrich EWANE}
\author[uy1]{Damaris Belle M. FOTSO}
\author[uy1]{Ren\'e NDOUNDAM}
\affiliation[uy1]{organization={University of Yaound\'e I},
 city={Yaound\'e}, country={Cameroon}}
\begin{abstract}
We study the Fisher--Rao geometry of parametric generative
transport and distinguish geometric path costs from observable
steganographic security. For regular statistical models,
the accumulated KL divergence between neighboring distributions
has a leading term proportional to Fisher path energy; a
constant-speed minimizing geodesic optimizes this leading term,
without changing the KL divergence between fixed endpoints.
We formulate natural-gradient optimization with explicit
regularity and feasibility requirements, and derive local
links to likelihood-ratio testing for exponential families.
To connect geometry to an actual observation model, we analyze a
keyed Gaussian-location signaling protocol with an encoder,
a decoder, an innocent-output distribution and an adversary
without the secret key. We obtain an exact Gaussian decoding
error and a nonasymptotic KL-security upper bound expressed
through per-stage squared Fisher displacements.
A constant-Fisher-speed schedule minimizes a certified
security upper bound at fixed receiver reliability.
Importantly, this does not by itself imply that a minimizing
geodesic outperforms every other constant-speed path.
Low-dimensional Gaussian simulations and a learned
VAE--MNIST latent-path experiment illustrate the modeling
distinctions. Operational distribution-preserving steganography
based on iMEC and Discop is discussed separately, without claiming
a matched numerical advantage for our procedure. The results
provide a precise starting point for observation-aware
Fisher-geometric steganographic design while leaving open
high-dimensional matched security and payload benchmarks.
\end{abstract}
\begin{keyword}
Fisher--Rao geometry \sep natural gradient \sep
generative steganography \sep Gaussian signaling \sep
relative entropy \sep steganalysis \sep keyed randomization
\end{keyword}
\end{frontmatter}


\section{Introduction}
Steganography aims to communicate information while avoiding
detection of the communication itself.
Security depends on the probability distribution of what an
adversary \emph{actually observes}, not on the geometry of an
unobserved parameter path alone. Distribution-preserving
generative techniques such as minimum-entropy-coupling methods
\cite{schroeder2023}, distribution-copy sampling
\cite{ding2023}, and sparse sampling \cite{wang2025}
construct operational encoding and decoding mechanisms.
Their role is distinct from the Fisher--Rao optimization of
a parametric statistical family.


Let $\{P_\theta:\theta\in\Theta\}$ be a regular model.
Its Fisher information metric provides an invariant local
notion of statistical displacement. We first describe the
path-energy limit of accumulated neighboring KL divergences
and the conditional role of natural gradients. We then
specify a secret-key signaling model and distinguish
message recovery from detectability of an \emph{innocent}
reference law. This separation is necessary: a reduction in
geodesic energy cannot be equated with a reduction in
steganographic detection error unless the observation
distributions are shown to depend on the chosen path.


The paper makes the following bounded contributions:
(i) a careful local KL/geodesic formulation, including its
endpoint invariance; (ii) a consistent natural-gradient and
curvature analysis; (iii) an explicit Gaussian secret-key
communication model with calculable receiver error and
adversarial KL bound; and (iv) a reproducibility-oriented
interpretation of synthetic and learned-generator diagnostics.
The Gaussian coding result is an analytically tractable
example of an observation-aware model, not an assertion of
new general-purpose steganographic capacity.


\section{Statistical model and Fisher geometry}
Let $p_\theta(x)>0$ be a smooth identifiable family on a
common dominating measure space, with regularity permitting
differentiation under the integral and finite score moments.
Define
\begin{equation}
 g_{ij}(\theta)=\E_\theta[
 \partial_i\log p_\theta(X)\,
 \partial_j\log p_\theta(X)].
 \label{eq:fisher}
\end{equation}
Assume $g(\theta)$ is positive definite on the regular
region of interest. If $\gamma:[0,1]\to\Theta$ is smooth,
let
\begin{align}
 L_g(\gamma)&=\int_0^1\|\dot\gamma(t)\|_{g(\gamma(t))}\,dt,\\
 E_g(\gamma)&=\int_0^1\|\dot\gamma(t)\|_{g(\gamma(t))}^{2}\,dt.
 \label{eq:energy}
\end{align}
In particular, $E_g(\gamma)\ge
L_g(\gamma)^2\ge d_g(\gamma(0),\gamma(1))^2$;
equality holds for constant-speed length-minimizing
geodesics when they exist.


\begin{proposition}[Sequential-KL small-step limit]
\label{prop:klpath}
For a $C^3$ path with image contained in a compact regular
parameter region and sufficiently uniform Taylor bounds,
let $t_k=k/N$ and
\[
S_N(\gamma)=\sum_{k=0}^{N-1}
 \KL(P_{\gamma(t_{k+1})}\Vert P_{\gamma(t_k)}).
\]
Then
\[
 S_N(\gamma)=\frac{1}{2N}E_g(\gamma)+O(N^{-2}).
\]
\end{proposition}
\begin{proof}
By the KL--Fisher local Taylor expansion,
$\KL(P_{\theta+\delta}\Vert P_\theta)
=\frac12\delta^\top g(\theta)\delta
+O(\|\delta\|^3)$.
Since $\delta_k=\dot\gamma(t_k)/N+O(N^{-2})$,
the $k$th term equals
$\|\dot\gamma(t_k)\|_g^2/(2N^2)+O(N^{-3})$.
Summation and Riemann quadrature prove the claim.
\end{proof}


\begin{remark}[Endpoint invariance]
\label{rem:endpoint}
For any two paths joining the same $\theta_0,\theta_1$,
$\KL(P_{\theta_1}\Vert P_{\theta_0})$ is identical.
Proposition~\ref{prop:klpath} is a result about a
\emph{sequence of neighboring distributions}; it is not an
unconditional security statement about a final sample.
Moreover $S_N(\gamma)\to0$ for fixed $\gamma$ as $N\to\infty$;
its invariant leading-order quantity is $NS_N$.
\end{remark}


\subsection{Exponential-family expansions}
For
$p_\theta(x)=h(x)\exp\{\theta^\top T(x)-\psi(\theta)\}$
in natural coordinates, let $g_{ij}=\psi_{ij}$ and
$T_{ijk}=\psi_{ijk}$.
For a small increment $\delta$,
\begin{align}
\KL(P_{\theta+\delta}\Vert P_\theta)
 &=\frac12g_{ij}\delta^i\delta^j+
   \frac13T_{ijk}\delta^i\delta^j\delta^k+
   O(\|\delta\|^4),\\
\KL(P_\theta\Vert P_{\theta+\delta})
 &=\frac12g_{ij}\delta^i\delta^j+
   \frac16T_{ijk}\delta^i\delta^j\delta^k+
   O(\|\delta\|^4),\\
\frac12d_g(\theta,\theta+\delta)^2
 &=\frac12g_{ij}\delta^i\delta^j+
   \frac14T_{ijk}\delta^i\delta^j\delta^k+
   O(\|\delta\|^4).
\end{align}
Consequently the forward and reverse KL divergences differ
from $\frac12d_g^2$ by $\pm T_{ijk}\delta^i\delta^j\delta^k/12$
to third order, respectively. The expansions require the
specified natural-coordinate and KL-orientation conventions.


\subsection{Gaussian reference geometry}
For the univariate normal $P_{\mu,\sigma}=\cN(\mu,\sigma^2)$,
\[
 g=\frac{d\mu^2+2\,d\sigma^2}{\sigma^2}
   =2\frac{d\mu^2+dy^2}{y^2},\qquad y=\sqrt2\sigma.
\]
This particular manifold is a complete hyperbolic plane with
sectional curvature $-1/2$. General exponential-family or
mixture manifolds do not inherit this completeness or
global uniqueness without separate assumptions.


\section{Optimization, curvature and statistical testing}
For a differentiable objective $f$ on a regular Fisher
manifold, the natural gradient is
$\operatorname{grad}_g f=g^{-1}\nabla f$.
An exact exponential-map descent step
$\theta^+=\Exp_\theta[-L^{-1}\operatorname{grad}_g f]$
satisfies
\[
 f(\theta^+)-f(\theta^\star)\le
 (1-\mu/L)\bigl(f(\theta)-f(\theta^\star)\bigr)
\]
whenever $f$ is $L$-Riemannian smooth and
$\mu$-geodesically strongly convex on an invariant
geodesically convex region containing the steps and optimum.
This guarantee does not automatically extend to an
additive coordinate step or covariance eigenvalue clipping.


For a full Gaussian $\cN(\bm\mu,\Sigma)$ with independent
coordinates $(\bm\mu,\operatorname{vech}\Sigma)$, the
covariance Fisher block is
\[
 (F_\Sigma)_{ab}=\frac12\operatorname{tr}
 (\Sigma^{-1}E_a\Sigma^{-1}E_b),
 \quad E_a=\partial\Sigma/\partial\theta_a,
\]
and the coordinate derivative of a scalar function is
$\partial f/\partial\theta_a=
\operatorname{tr}[(\nabla_\Sigma f)^\top E_a]$.
Off-diagonal independent covariance coordinates contribute
twice to the corresponding matrix derivative.
For $\Sigma=I_2$ this block is
$\operatorname{diag}(1/2,1,1/2)$ in the ordering
$(\sigma_{11},\sigma_{12},\sigma_{22})$.


The Jacobi equation
$J''+R(J,\dot\gamma)\dot\gamma=0$
describes infinitesimal geodesic sensitivity. In constant
sectional curvature $K=-a^2$ a unit transverse Jacobi factor
equals $\sinh(at)/a$; negative scalar curvature in dimension
greater than two does not determine the sign of every
sectional curvature. Geometric ball volume, packing, and
operational communication capacity are distinct notions.


Under simple hypotheses $H_0=P_{\theta_c}$ (cover) and
$H_1=P_{\theta_s}$ (stego) in a regular exponential family,
the normalized log-likelihood ratio is
\[
 \Lambda_n=(\theta_s-\theta_c)^\top
  (\widehat\eta_n-\eta_c)-
  \KL(P_{\theta_c}\Vert P_{\theta_s}).
\]
At fixed type-I level under $H_0$, the optimal asymptotic
type-II error exponent is
$\KL(P_{\theta_c}\Vert P_{\theta_s})$ by the
Chernoff--Stein lemma. This result does not alone imply
a finite-sample safe Fisher-radius guarantee.




\section{An observation-aware keyed Gaussian signaling model}
\label{sec:keyed}
We now give a concrete communication model in which geometry,
decoding reliability and adversarial indistinguishability have
separately identifiable mathematical quantities.
The model is an illustrative keyed Gaussian channel, not an
image steganographic system and not a claim that the following
one-bit spread-spectrum construction is new.


\subsection{Alice, Bob, Eve and the innocent distribution}
Consider a single uniform message bit
$M\in\{-1,+1\}$ and a secret key
$K=(K_1,\ldots,K_N)$ of independent uniform
Rademacher signs, shared by Alice and Bob but not Eve.
Let $\sigma>0$ and let $a_1,\ldots,a_N$ be public
nonnegative amplitudes.
Alice transmits
\begin{equation}
Y_k= M K_k a_k+\varepsilon_k,\qquad
\varepsilon_k\overset{\mathrm{iid}}{\sim}\cN(0,\sigma^2).
\label{eq:encoder}
\end{equation}
In the innocent condition $H_0$, Eve instead observes
$Y_k=\varepsilon_k$ independently, so
$Q_0=\bigotimes_{k=1}^N\cN(0,\sigma^2)$.
In the active condition $H_1$, Eve knows the amplitude
schedule and protocol but not the secret key or bit.
Averaging over the key gives
\begin{equation}
Q_1=\bigotimes_{k=1}^N
\left\{\tfrac12\cN(a_k,\sigma^2)+
       \tfrac12\cN(-a_k,\sigma^2)\right\}.
\label{eq:cover-stego}
\end{equation}
The independent secret signs are essential: a fixed,
public or reused key need not provide the product mixture
in Eq.~\eqref{eq:cover-stego}.
This observer distinguishes innocent noise
from an active secret-key-modulated Gaussian signal;
the decoder has an additional secret key. The example
therefore supplies a fully specified threat model, not
the distribution preservation of a natural image generator.


\subsection{Exact decoding and finite-sample KL control}
\begin{theorem}[Receiver reliability and Eve's KL bound]
\label{thm:gaussian_security}
Under Eqs.~\eqref{eq:encoder}--\eqref{eq:cover-stego},
assume $M$ and $K$ are independent and uniform and Bob
receives $Y_1,\ldots,Y_N$ without additional channel
distortion. Put
$\rho_k=a_k^2/\sigma^2$ and
$\rho=\sum_{k=1}^N\rho_k$.
If $\rho>0$, the optimum equally likely binary-message
decoder with the key is
\begin{equation}
\widehat M=\sgn\left(\sum_{k=1}^N K_k a_kY_k\right),
\end{equation}
with exact error probability
\begin{equation}
\Pr\{\widehat M\ne M\}=\PhiN(-\sqrt\rho).
\label{eq:pe}
\end{equation}
For an adversary without the key, the active-versus-innocent
divergence satisfies
\begin{equation}
0\le \KL(Q_1\Vert Q_0)
\le\sum_{k=1}^N\log\cosh(\rho_k)
\le\tfrac12\sum_{k=1}^N\rho_k^2.
\label{eq:quartic-security}
\end{equation}
The total-variation advantage of every binary
active-versus-innocent detector is bounded by
\begin{equation}
\TV(Q_1,Q_0)
\le \sqrt{\tfrac12\KL(Q_1\Vert Q_0)}
\le \tfrac12\sqrt{\sum_k\rho_k^2}.
\label{eq:pinsker}
\end{equation}
\end{theorem}
\begin{proof}
Conditionally on $(M,K)$, the weighted correlation
$Z=\sum_k K_ka_kY_k$ is normal with mean
$M\sum_ka_k^2$ and variance
$\sigma^2\sum_ka_k^2$. Its sign is the
likelihood-ratio decoder and Eq.~\eqref{eq:pe} follows.
For one coordinate, relative to
$q_0(y)=\varphi_\sigma(y)$,
the mixture likelihood ratio is
\[
\frac{q_1(y)}{q_0(y)}
=e^{-\rho_k/2}
\cosh(a_ky/\sigma^2).
\]
A direct Gaussian moment calculation yields
\[
\chi^2(q_1,q_0)
=\E_{q_0}\!\left[
\left(q_1/q_0-1\right)^2\right]
=\cosh(\rho_k)-1.
\]
Jensen's inequality gives
$\KL(q_1\Vert q_0)
=\E_{q_1}\log(q_1/q_0)
\le\log(1+\chi^2(q_1,q_0))
=\log\cosh(\rho_k)$.
The elementary bound
$\log\cosh x\le x^2/2$ and product additivity
of KL give Eq.~\eqref{eq:quartic-security}.
Pinsker's inequality gives Eq.~\eqref{eq:pinsker}.
\end{proof}


\begin{corollary}[Uniform allocation minimizes the certified
security upper bound]
\label{cor:uniform}
For fixed $N$, $\sigma$ and receiver squared
signal-to-noise ratio $\rho=\sum_k\rho_k$,
the upper bound $\tfrac12\sum_k\rho_k^2$ is
minimized over $\rho_k\ge0$ by the uniform
allocation $\rho_k=\rho/N$, yielding
\begin{equation}
\KL(Q_1\Vert Q_0)\le\frac{\rho^2}{2N},
\qquad
\TV(Q_1,Q_0)\le\frac{\rho}{2\sqrt N},
\qquad
\Pr(\widehat M\ne M)=\PhiN(-\sqrt\rho).
\label{eq:allocation}
\end{equation}
\end{corollary}
\begin{proof}
By Cauchy--Schwarz,
$\sum_k\rho_k^2\ge
(\sum_k\rho_k)^2/N$, with equality
at uniform $\rho_k$.
The conclusion optimizes the stated \emph{certified
upper bound}. It does not establish that uniform
allocation uniquely minimizes the \emph{exact} KL,
and it is not a claim of a universal optimal
Fisher geodesic.
\end{proof}


\begin{remark}[Where Fisher geometry enters]
\label{rem:geometrysecurity}
For the constant-variance Gaussian location family
$\{\cN(\mu,\sigma^2)\}$, the Fisher line element is
$d\mu^2/\sigma^2$. Hence each $\rho_k$ is precisely
the squared Fisher displacement for the elementary
mean perturbation $0\mapsto a_k$.
The security certificate in
Eq.~\eqref{eq:quartic-security} is
\emph{quartic} in the local Fisher displacement,
rather than proportional to the quadratic
neighboring-KL energy
$S_N$ of Proposition~\ref{prop:klpath}.
This distinction is structural: averaging over
secret signs cancels the first-order shift in the
adversary's observation law.
Equal Fisher squared displacement per stage
minimizes the stated bound under a fixed
receiver reliability. This is a
\emph{resource-allocation} statement in a
flat location model, not proof that an arbitrary
geodesic minimizes security leakage.
\end{remark}


\begin{remark}[Boundaries of the result]
The adversarial security guarantee concerns the
\emph{average over fresh uniformly random secret keys}.
If a key is disclosed, reused across correlated
transmissions or chosen nonuniformly, the
observer's law may differ. The theorem concerns
one recoverable bit and a Gaussian location
model. Increasing $N$ while allocating
$\rho_k=\rho/N$ holds Bob's total SNR fixed
but consumes more channel uses and fresh key bits.
The result does not constitute constant-rate
covert communication or a positive asymptotic
payload rate. It is not a claim of security for
MNIST, CIFAR-10, iMEC or Discop.
\end{remark}


\subsection{Consequence for geometric transport claims}
For a fixed pair of terminal distributions, neither
the exact decoder error nor the observation
divergence of an actual channel follows from the
endpoint-independent quantity $S_N$ without
a channel/observation construction.
The keyed model proves a different, appropriately
specified implication: a secret-key mixture makes
the relative entropy of Eve's observation law
bounded by fourth-order Fisher displacements
at fixed receiver reliability.
In general statistical manifolds a comparable
result would require analyzing the
key-marginalized observation density and
its higher-order score tensors, not merely
invoking a minimizing geodesic.


\section{Numerical illustrations and reproducibility}
\label{sec:experiments}
\subsection{Analytic reference experiments}
Low-dimensional Gaussian and Gaussian-mixture
illustrations evaluate geodesic integration, the
coordinate-consistent Fisher metric and curvature
diagnostics. Numerical convergence claims must
be evaluated on the corresponding corrected
script outputs; unverified percentages or
runtime factors are not used as results here.
For the Gaussian location model of
Theorem~\ref{thm:gaussian_security}, the
testable quantities are (i) decoding error
$\PhiN(-\sqrt\rho)$, (ii) the active/innocent
likelihood-ratio distribution, (iii) exact or
numerically integrated KL versus
$\frac12\sum_k\rho_k^2$, and (iv) the
resource-allocation curve $\rho_k=\rho/N$.
No empirical values are asserted for tests
not represented by archived outputs.


\subsection{Learned-generator latent paths}
We also consider an exploratory VAE--MNIST
simulation with a two-dimensional latent
variable $z$ and Bernoulli decoder
$p_\phi(x\mid z)$. Its pullback Fisher
metric can be computed as
\[
F(z)=J_\pi(z)^\top\!
\operatorname{diag}\left(
\frac1{\pi_i(z)(1-\pi_i(z))}
\right)J_\pi(z).
\]
The available code compares a straight latent
interpolation with a path obtained by
Fisher-preconditioned descent toward a
target code. The latter is not in general a
Levi--Civita minimizing geodesic.
The code's detector is trained to distinguish
samples near source codes from images generated
along straight interpolations, which differs
from the innocent/no-message protocol of
Theorem~\ref{thm:gaussian_security}.
The repository contains outputs for three
digit-class pairs, but these results are
exploratory and do not measure message-decoding
accuracy or establish a matched Fisher--Rao
security advantage. An additional matched
evaluation would be required.


\subsection{Operational references}
Minimum-entropy coupling and
distribution-copy embedding provide
operational methods for steganographic
generation \cite{schroeder2023,ding2023}.
They are not Euclidean parameter-path
optimizers. Independent reproduction
and detector-health experiments can be
used as protocol references, but their
detector AUCs are not attributed to the
Gaussian signaling model or the
VAE latent-path illustration.
No additional iMEC/Discop generation
is necessary to establish the analytic
statements of Section~\ref{sec:keyed}.


\section{Discussion and limitations}
The main conceptual distinction is between
optimization geometry and observable
steganographic security. The geodesic
energy principle is standard and
characterizes a leading-order mathematical
path cost. A fixed terminal law is
not made harder to distinguish merely
because it is reached by a geodesic.


The explicit keyed Gaussian protocol
specifies the innocent distribution,
the secret-key-marginalized active
distribution, Bob's optimal decoder and
a finite-sample KL certificate for Eve.
Its bound captures cancellation due to
key averaging, not a general capacity
law for generative media. The
uniform-energy allocation result is
limited to a fixed total receiver SNR
and a fixed number of stages.
The claimed contribution is therefore
the clear observation-aware
specialization and connection to
Fisher displacement, rather than
priority over existing spread-spectrum
or covert-Gaussian-channel constructions.


The available learned-generator
experiment does not operate the
keyed signaling protocol, lacks an
explicit message decoder and does not
prove that the natural-gradient path
minimizes Fisher energy. Consequently
the Gaussian theorem is not validated
as an image-steganography security
result. A future method that encodes
a message by a route-dependent generator
would require an explicit receiver,
a key and a joint innocent/active
observation law and must satisfy
payload-recovery and detection
criteria on the same data. Likewise,
high-dimensional exact-Fisher
computational costs remain outside
our demonstrated scope.


\section{Conclusion}
Fisher--Rao geometry offers an
intrinsic framework for analyzing
parameter transport, but geometric
optimality alone is not an operational
steganographic security proof.
We establish the small-step
sequential-KL/energy relation
under regularity conditions and
preserve the distinctions between
endpoints, paths and observations.
For a fully specified keyed Gaussian
signaling model, we derive exact
receiver reliability and a nonasymptotic
upper bound on the relative entropy of
the adversary's observations.
A constant-Fisher-energy-per-stage
schedule minimizes that certificate
at fixed decoding SNR. The result
is an analytically transparent
benchmark for observation-aware
Fisher-geometric design, rather than
a demonstrated universal improvement
over current operational generative
steganography.


\begin{thebibliography}{99}
\bibitem{cachin2004}
C. Cachin, An information-theoretic model for steganography,
Information and Computation 192 (2004) 41--56.


\bibitem{schroeder2023}
C. Schroeder de Witt, S. Sokota, J. Z. Kolter, J. Foerster,
M. Strohmeier, Perfectly Secure Steganography Using
Minimum Entropy Coupling, ICLR, 2023.


\bibitem{ding2023}
J. Ding, K. Chen, Y. Wang, N. Zhao, W. Zhang, N. Yu,
Discop: Provably Secure Steganography in Practice
Based on Distribution Copies, IEEE Symposium on
Security and Privacy, 2023.
doi:10.1109/SP46215.2023.10179287.


\bibitem{wang2025}
Y. Wang et al., SparSamp: Efficient Provably Secure
Steganography Based on Sparse Sampling,
USENIX Security Symposium, 2025.


\bibitem{amari2000}
S.-i. Amari, H. Nagaoka,
Methods of Information Geometry,
American Mathematical Society, 2000.


\bibitem{amari1998}
S.-i. Amari, Natural gradient works efficiently
in learning, Neural Computation 10 (1998) 251--276.


\bibitem{ollivier2017}
Y. Ollivier, L. Arnold, A. Auger, N. Hansen,
Information-Geometric Optimization Algorithms:
A Unifying Picture via Invariance Principles,
Journal of Machine Learning Research 18 (2017) 1--65.


\bibitem{boumal2019}
N. Boumal, P.-A. Absil, C. Cartis,
Global rates of convergence for nonconvex
optimization on manifolds,
IMA Journal of Numerical Analysis 39 (2019) 1--33.


\bibitem{cover2006}
T. Cover, J. Thomas, Elements of Information
Theory, 2nd ed., Wiley, 2006.


\bibitem{fridrich2009}
J. Fridrich, Steganography in Digital Media:
Principles, Algorithms, and Applications,
Cambridge University Press, 2009.


\bibitem{rygaard2026}
F. M. Rygaard, S. Hauberg,
GEORCE: A fast new control algorithm
for computing geodesics,
Information Geometry, 2026.
doi:10.1007/s41884-026-00192-9.
\end{thebibliography}
\end{document}